\section{Đánh giá triết gia và các nhà kinh tế}

\subsection{Quan điểm của Lê-nin}
\subsubsection{Lý luận}
Sau khi biết được về lý thuyết giá trị thặng dư và quy luật bóc lột giá trị mà người lao động tạo ra trong quá trình sản xuất. Lê-nin đã phát triển lý thuyết của Marx trong bối cảnh chủ nghĩa tư bản độc quyền và chủ nghĩa đế quốc dần phát triển. Với Lê nin: "Chủ nghĩa đế quốc là giai đoạn tốt cùng của chủ nghĩa tư bản" \cite{LeninImperial} cũng chính là tựa đề cho quyển sách của ông. Theo Lênin, trong giai đoạn CNTB tự do cạnh tranh, các công ty nhỏ có thể tồn tại và phát triển. Tuy nhiên, khi CNTB phát triển, các doanh nghiệp lớn dần thống trị thị trường thông qua sự cạnh tranh khốc liệt. Kết quả là tư bản và sản xuất ngày càng tập trung vào tay một số ít doanh nghiệp độc quyền. Điều này cho phép các tập đoàn tư bản kiểm soát thị trường và gia tăng việc chiếm đoạt giá trị thặng dư từ lao động. Đây là đặc điểm cốt lõi của Chủ nghĩa Đế quốc (CNĐQ): Sự thống trị của các tổ chức độc quyền trên thị trường toàn cầu, gây ra sự phân cực giàu nghèo và bất bình đẳng kinh tế. \cite{ch1_Lenin_imper}

Trong giai đoạn đầu của CNTB, các nước phát triển tập trung vào việc xuất khẩu hàng hóa. Nhưng đến giai đoạn đế quốc chủ nghĩa, các quốc gia tư bản phát triển chuyển sang xuất khẩu tư bản (vốn), tức là đầu tư trực tiếp vào các nước lạc hậu hoặc đang phát triển để thu lợi nhuận. Điều này làm gia tăng sự phụ thuộc kinh tế của các nước kém phát triển vào các nước tư bản phát triển, và là một hình thức bóc lột mới. Từ đây các nước tư bản phát triển bắt đầu tranh giành thuộc địa và tài nguyên thiên nhiên ở các nước kém phát triển. Lênin lập luận rằng, việc phân chia thuộc địa và ảnh hưởng chính trị giữa các cường quốc tư bản là biểu hiện rõ ràng của CNĐQ. Những nước này không chỉ tìm kiếm thị trường mới mà còn khai thác tài nguyên, vốn lao động giá rẻ để bóc lột đến tốt cùng giá trị thặng dư từ các quốc gia đó. \cite{ch4_Lenin_imper}

Khi CNTB phát triển lên giai đoạn đế quốc, các mâu thuẫn xã hội và kinh tế ngày càng gay gắt, đặc biệt là mâu thuẫn giữa lực lượng sản xuất phát triển cao và quan hệ sản xuất TBCN bó hẹp. Mâu thuẫn giữa các cường quốc đế quốc về việc phân chia tài nguyên và thị trường toàn cầu dẫn đến các cuộc xung đột lớn, như các cuộc chiến tranh thế giới. Qua phân tích đó, Lenin đã làm rõ rằng học thuyết giá trị thặng dư không chỉ giải thích cơ chế bóc lột trong sản xuất tư bản chủ nghĩa mà còn giúp lý giải các hiện tượng kinh tế – chính trị của chủ nghĩa đế quốc, nhờ đó khẳng định tính thời sự và khả năng phát triển của lý luận Marx trong việc phân tích chủ nghĩa tư bản ở giai đoạn mới.

\subsubsection{Thực tiễn}

Tuy nhiên câu hỏi mới đặt ra, nếu không còn tư bản tư nhân thì giá trị thặng dư sẽ đi đâu? Với Lê-nin, giá trị thặng dư bị nhà tư bản chiếm đoạt và bóc lột từ sức lao động của người dân thì ở nhà nước xã hội chủ nghĩa, giá trị thặng dư vẫn được tạo ra trong quá trình sản xuất nhưng nó thuộc về nhà nước - đại diện của giai cấp công nhân. Để từ đó phát triển công nghiệp hóa, kinh tế, phúc lợi xã hội cho đất nước - ở đây là phục vụ cho toàn dân.\cite{tax-in-kind} Điều này có nghĩa rằng, bản thân giá trị thặng dư vẫn tồn tại, nhưng cách phân phối thay đổi và trong quá trình quá độ lên XHCN, ta vẫn chấp nhận tạm thời hình thức kinh tế thị trường.

Thực tế thì sau khi thành lập Liên Xô, Lê-nin đã đưa ra \textbf{"chính sách kinh tế mới - NEP"} với quan điểm vẫn cho phép một mức độ thị trường và tư bản nhỏ trong khi nhà nước sẽ kiểm soát các ngành kinh tế lớn, quan trọng và trọng điểm của quốc gia. Điều này khiến giá trị thặng dư vẫn được tạo ra trong khu vực tư nhân nhưng nhà nước Xô viết thu lại được thông qua thuế và kiểm soát kinh tế. Lê-nin đã từng nói rằng trong giai đoạn chuyển tiếp: "nhà nước công nhân có thể sử dụng \textbf{tư bản nhà nước} để phát triển lực lượng sản xuất".\cite{State-and-revol}. Chung quy lại, nhà nước Xô viết không xóa bỏ ngay toàn bộ quan hệ kinh tế cũ mà kiểm soát và hướng giá trị thặng dư phục vụ xã hội. Giá trị thặng dư không biến mất khi cách mạng thành công, vấn đề là ai kiểm soát và phân phối nó. Ở đây nhà nước Xã hội chủ nghĩa là nhà nước đại diện cho giai cấp công nhân quản lý lấy.

\subsection{Quan điểm phê phán: các nhà tư bản, kinh tế học}
Bên cạnh những quan điểm ủng hộ, học thuyết giá trị thặng dư của Karl Marx cũng vấp phải nhiều phê phán từ các nhà kinh tế học thuộc trường phái kinh tế tư bản và kinh tế học chính thống. Một trong những phê phán nổi bật đến từ Eugen von Böhm-Bawerk, người cho rằng lý thuyết giá trị lao động của Marx không thể giải thích đầy đủ sự hình thành giá cả trong nền kinh tế thị trường.\cite{EugenCritic} Theo ông, Marx đã bỏ qua vai trò của thời gian, vốn và sở thích tiêu dùng trong việc xác định giá trị hàng hóa. Ngoài ra, Böhm-Bawerk cũng chỉ ra cái gọi là “mâu thuẫn nội tại” trong tác phẩm Capital, khi Marx vừa khẳng định giá trị do lao động quyết định, vừa thừa nhận rằng giá cả thị trường có thể khác biệt do sự phân bổ tư bản giữa các ngành. Các nhà kinh tế học tân cổ điển sau này như Paul Samuelson tiếp tục phát triển các phê phán này, cho rằng lý thuyết giá trị thặng dư không cần thiết để giải thích lợi nhuận, bởi trong kinh tế học hiện đại, giá cả và thu nhập của các yếu tố sản xuất được xác định thông qua cung – cầu và năng suất cận biên. Từ góc nhìn đó, lợi nhuận của nhà tư bản không phải là kết quả của sự chiếm đoạt giá trị thặng dư từ lao động, mà được xem như phần thưởng cho việc đầu tư vốn, chấp nhận rủi ro và tổ chức sản xuất. Những phê phán này cho thấy sự khác biệt căn bản giữa kinh tế học Marxist và kinh tế học chính thống trong cách lý giải nguồn gốc của giá trị và lợi nhuận trong nền kinh tế thị trường.

\subsubsection{Eugen von Böhm-Bawerk}
Một trong những nhà kinh tế học tư bản tiêu biểu phê phán học thuyết giá trị thặng dư của Karl Marx là Eugen von Böhm-Bawerk, đại diện của trường phái kinh tế học Áo. Trong tác phẩm Karl Marx and the Close of His System,\cite{EugenCritic} Böhm-Bawerk cho rằng lý thuyết giá trị lao động của Marx tồn tại những mâu thuẫn nội tại và không thể giải thích đầy đủ cơ chế hình thành giá cả trong nền kinh tế tư bản chủ nghĩa. Theo Marx, giá trị của hàng hóa được quyết định bởi lượng lao động xã hội cần thiết để sản xuất ra nó, và lợi nhuận của nhà tư bản thực chất là giá trị thặng dư được tạo ra từ lao động của công nhân.\cite{Capital} Tuy nhiên, Böhm-Bawerk lập luận rằng trong thực tế, giá cả thị trường không chỉ phụ thuộc vào lao động mà còn chịu ảnh hưởng của nhiều yếu tố khác như mức độ khan hiếm của hàng hóa, nhu cầu của người tiêu dùng, cũng như vai trò của vốn và thời gian trong quá trình sản xuất.

Böhm-Bawerk đặc biệt nhấn mạnh đến cái mà ông gọi là “vấn đề chuyển hóa” trong hệ thống lý luận của Marx. Theo ông, trong tập I của tác phẩm Capital\cite{Capital}, Marx khẳng định rằng giá trị hàng hóa được quyết định bởi lao động, nhưng đến các phân tích sau đó lại thừa nhận rằng lợi nhuận giữa các ngành có xu hướng ngang bằng nhau do sự di chuyển của tư bản. Điều này dẫn đến việc giá cả thực tế của hàng hóa không hoàn toàn phản ánh lượng lao động chứa đựng trong chúng. Từ đó, Böhm-Bawerk cho rằng Marx đã không giải thích thỏa đáng sự chuyển hóa từ giá trị sang giá cả sản xuất, khiến cho hệ thống lý luận về giá trị thặng dư trở nên thiếu nhất quán. Ngoài ra, Böhm-Bawerk còn cho rằng lợi nhuận của nhà tư bản không nhất thiết phải được hiểu như sự bóc lột lao động, mà có thể được giải thích như phần bù đắp cho việc đầu tư vốn, chấp nhận rủi ro và trì hoãn tiêu dùng trong quá trình sản xuất.\cite{EugenCritic}

Những lập luận của Böhm-Bawerk đã trở thành một trong những phê phán có ảnh hưởng lớn đối với kinh tế học Marxist và mở ra nhiều cuộc tranh luận kéo dài trong kinh tế học về nguồn gốc của giá trị và lợi nhuận trong nền kinh tế thị trường. Mặc dù những phê phán của Eugen von Böhm-Bawerk đối với học thuyết giá trị thặng dư của Karl Marx đã có ảnh hưởng lớn trong kinh tế học phương Tây, đặc biệt là trong việc chỉ ra những vấn đề liên quan đến sự chuyển hóa từ giá trị sang giá cả, song các nhà kinh tế học Marxist sau này cũng đưa ra nhiều phản biện nhằm bảo vệ và phát triển hệ thống lý luận của Marx. Một số học giả cho rằng những điểm mà Böhm-Bawerk gọi là “mâu thuẫn” thực chất xuất phát từ việc ông diễn giải chưa đầy đủ phương pháp phân tích của Marx trong tác phẩm Capital, nơi Marx phân biệt giữa khái niệm giá trị ở mức độ trừu tượng và sự vận động của giá cả trong thực tiễn thị trường. Ngoài ra, nhiều nghiên cứu sau này cũng cố gắng giải quyết vấn đề chuyển hóa bằng các mô hình toán học và cách tiếp cận mới trong kinh tế học Marxist như Maurice Dobb, Andrew Kliman,... Vì vậy, thay vì làm suy yếu hoàn toàn học thuyết giá trị thặng dư, các phê phán của Böhm-Bawerk đã góp phần thúc đẩy những cuộc tranh luận học thuật sâu rộng về bản chất của giá trị, lợi nhuận và quan hệ sản xuất trong chủ nghĩa tư bản. Qua đó có thể thấy rằng học thuyết giá trị thặng dư vẫn giữ vai trò quan trọng như một công cụ lý luận để phân tích những đặc điểm và mâu thuẫn của nền kinh tế tư bản chủ nghĩa.

\section{Kết luận và Ý nghĩa thực tiễn đối với Việt Nam}

\subsection{Nhận diện đúng bản chất}

Mặc dù có những thay đổi đáng kể như thế, nhưng sự thay đổi biểu hiện bên ngoài không tạo ra sự thay đổi về mặt bản chất của chủ nghĩa tư bản. Quan hệ kinh tế của chủ nghĩa tư bản vẫn dựa trên cơ sở là hình thức sở hữu tư nhân về tư liệu sản xuất. Cán cân quyền lực kinh tế vẫn nghiêng về giai cấp tư sản \cite{tapchi_congthuong_gttd}.

Luận điểm của C. Mác và Ph. Ăng-ghen về chế độ người bóc lột người trong xã hội tư bản vẫn giữ nguyên giá trị khoa học. Mặc dù đã có sự thay đổi và điều chỉnh, song có thể khẳng định: bản chất bóc lột của chủ nghĩa tư bản không thay đổi \cite{tran_hieu_gttd}. Học thuyết giá trị thặng dư với tư cách là một hệ thống lý luận phong phú và sâu sắc về kinh tế thị trường nhằm vận dụng vào công cuộc xây dựng và phát triển kinh tế -- xã hội trong nền kinh tế tri thức \cite{tapchi_congthuong_gttd}.

\subsection{Định hướng chiến lược cho Việt Nam}

Nghiên cứu học thuyết giá trị thặng dư có ý nghĩa to lớn đối với nước ta trong quá trình xây dựng và phát triển kinh tế theo hướng kinh tế tri thức. Ở Việt Nam trong thời kỳ quá độ lên chủ nghĩa xã hội hiện nay, học thuyết này có ý nghĩa hiện thực to lớn cho quá trình phát triển kinh tế hướng đến nền kinh tế tri thức. Cần vận dụng học thuyết một cách thông minh, sáng tạo nhưng đảm bảo tính khoa học phù hợp với điều kiện thực tiễn của Việt Nam \cite{tapchi_congthuong_gttd}.

Muốn tối ưu hóa lợi nhuận, Việt Nam cần thực hiện phương pháp sản xuất giá trị thặng dư siêu ngạch. Cần nhận thức rõ phương pháp đem lại giá trị thặng dư và lợi nhuận cao là tìm kiếm giá trị thặng dư nhờ áp dụng khoa học kỹ thuật, hiện đại hóa sản xuất, hạ giá thành sản phẩm \cite{tapchi_congthuong_gttd}.

Việt Nam cần ưu tiên phát triển nền kinh tế tri thức như một điều kiện sống còn trong thế giới đang vận động rất mau lẹ. Xác định rõ sự phát triển của khoa học, công nghệ là điều kiện cơ bản để hình thành và phát triển kinh tế tri thức \cite{tapchi_congthuong_gttd}.

Cần đầu tư hơn nữa cho khoa công nghệ, chú trọng công tác giáo dục, thực hiện chính sách thu hút người lao động có trình độ cao, tránh nguy cơ chảy máu chất xám,... Hoàn thiện hệ thống pháp luật về quyền sở hữu trí tuệ phù hợp với Việt Nam hiện nay, bảo vệ lợi ích của nước mua công nghệ, đồng thời khuyến khích sự sáng tạo từ trong nước \cite{tapchi_congthuong_gttd}.

Xuất phát điểm của Việt Nam thấp hơn so với các nước khác trong khu vực và trên thế giới, do đó Việt Nam cần thực hiện chiến dịch ``đi tắt, đón đầu'', học tập những thành tựu khoa học công nghệ và kinh nghiệm quản lý tiên tiến của các nước trên thế giới nhằm nâng cao hiệu quả hoạt động sản xuất -- kinh doanh \cite{tapchi_congthuong_gttd}.
